\documentclass[10pt,a4paper]{amsart}
\usepackage[margin=19mm]{geometry}
\usepackage{amsmath,amssymb,mathtools}
\usepackage[scr=rsfs,cal=euler]{mathalfa}
\usepackage{xfrac}
\usepackage[unicode,hidelinks,bookmarks=false]{hyperref}
\newcommand{\ie}{i.e.\ }
\newcommand{\cf}{cf.\ }
\newcommand{\resp}{respectively\ }
\newcommand{\Subsection}{\subsection}
\providecommand{\nobreakdash}{\nobreak\hbox{-}}
\setlength{\emergencystretch}{2em}
\multlinegap0pt
\usepackage{bm}
\usepackage{bbm}
\usepackage{dsfont}
\usepackage{xspace}
\usepackage{cancel}
\usepackage{mathtools}
\usepackage{amsthm}
\makeatletter
\@ifundefined{thm}{\newtheorem{thm}{Thm}}{}
\@ifundefined{lem}{\newtheorem{lem}[thm]{Lemma}}{}
\makeatother
\makeatletter
\newcommand{\cS}{\ensuremath{{\mathcal{S}}}}
\expandafter\ifx\csname customcor\endcsname\relax
\newenvironment{customcor}[1]
	{\renewcommand
	\theinnercustomcor{#1}\innercustomcor}
	{\endinnercustomcor}
\fi
\expandafter\ifx\csname customlem\endcsname\relax
\newenvironment{customlem}[1]
	{\renewcommand
	\theinnercustomlem{#1}\innercustomlem}
	{\endinnercustomlem}
\fi
\expandafter\ifx\csname customprop\endcsname\relax
\newenvironment{customprop}[1]
	{\renewcommand
	\theinnercustomprop{#1}\innercustomprop}
	{\endinnercustomprop}
\fi
\expandafter\ifx\csname customthm\endcsname\relax
\newenvironment{customthm}[1]
    {\renewcommand
    \theinnercustomthm{#1}\innercustomthm}{\endinnercustomthm}
\fi
\makeatother
\title[Exceptional sequences of type $B\_n/C\_n$ and those in the abe]{Exceptional sequences of type $B\_n/C\_n$ and those in the abelian tube}
\author{Kiyoshi Igusa, Emre Sen}
\date{Source: arXiv:2310.01700v3 (CC BY 4.0)}
\begin{document}
\maketitle

\noindent\emph{Source and licence.} Exceptional sequences of type $B_n/C_n$ and those in the abelian tube. Version arXiv:2310.01700v3 (CC BY 4.0), distributed under the Creative Commons Attribution 4.0 International licence, as recorded in the arXiv licence metadata for this exact version and shown on the abstract page https://arxiv.org/abs/2310.01700v3. Full rights notice: licenses/arxiv-2310.01700-CC-BY-4.0.txt.\par\smallskip

\noindent\textit{Editorial scope.} Counted unit: the Lem whose source label is \texttt{lem: second recursive count of S(mu,l)} in the article (source0.tex lines 2367--2372, source label \texttt{lem: second recursive count of S(mu,l)}), delivered together with the article's own complete original proof of it (source0.tex lines 2374--2399, one \texttt{proof} environment of 26 lines). No mathematics is written, repaired or re-derived: statement and proof are the article's own text line for line. Required context retained uncounted: eq: recursive count of S(mu,l) (lines 2359--2361). Every definition, display and label of those lines is retained unchanged. Omitted from this package and not redistributed: 1--2358, 2362--2366, 2373--2373, 2400--2592. Nothing inside a retained excerpt is omitted or altered: every retained body is delivered exactly as the article's lines are written, so no omission receipt of any kind accompanies this package. Citations: the retained excerpts carry no citation command at all, so no bibliography entry of the article is transcribed and no reference is redistributed. Reference adaptation: none was needed and none was made; every reference inside the delivered unit resolves inside it, so no unresolved reference or citation is produced. Printed slips kept literal: no printed word, symbol, hypothesis or number of the counted statement or of its proof was altered. Layout and class: the article's own class is \texttt{amsart} with packages including geometry, graphicx, amssymb, amsmath, color, pdfsync, tikz, tikz-cd, xy; the wrapper is the lane's pinned \texttt{amsart} editorial wrapper at 10pt on A4, which restates the theorem-like environments the retained text needs and adds the article's own macro definitions that the retained text needs (\texttt{\textbackslash cS}), with extra wrapper packages (bm, bbm, dsfont, xspace, cancel, mathtools). The delivered numbering is the unit's own. Assets: no retained span contains a figure environment, a graphics-inclusion command, a PDF-page inclusion, a table or a TikZ picture; no image file is copied into this package, the package declares no asset dependency and no image file is redistributed with it. Licence: version arXiv:2310.01700v3 of this article is distributed under the Creative Commons Attribution 4.0 International licence, as recorded in the arXiv licence metadata for this exact version and shown on the abstract page https://arxiv.org/abs/2310.01700v3. A full rights notice is delivered at licenses/arxiv-2310.01700-CC-BY-4.0.txt. Characters: every retained excerpt is pure ASCII, so the delivered engine, XeLaTeX with its default Latin Modern text font, reports no missing character.
\begin{equation}\label{eq: recursive count of S(mu,l)}
    (k+1)|\cS_n(\mu;\lambda)|=\sum_{\lambda'}|\cS_n(\mu';\lambda')|+\sum_{\lambda''}n_c''2X(a,b)|\cS_n(\mu;\lambda'')|
\end{equation}

\begin{lem}\label{lem: second recursive count of S(mu,l)}
With the same notation as in \eqref{eq: recursive count of S(mu,l)} we have:
\[
n|\cS_n(\mu;\lambda)|=\sum_{\lambda'} (\mu+a+1)|\cS_n(\mu';\lambda')|+\sum_{\lambda''} n_c''(c+1)X(a,b) |\cS_n(\mu;\lambda'')|.
\]
\end{lem}

\begin{proof}
We construct a mapping $\varphi:\widetilde\cS_n(\mu;\lambda)\to \widetilde\cS_n^p(\mu;\lambda)$. Each term on the right hand side of \eqref{eq: recursive count of S(mu,l)} corresponds to an element of $\widetilde\cS_n^p(\mu;\lambda)$. If we multiply this with the size of its inverse image in $\widetilde\cS_n(\mu;\lambda)$ we will obtain the required formula for the size of $\widetilde\cS_n(\mu;\lambda)$. 

For any $G\in \cS_n(\mu;\lambda)$ let $G_0,G_1,\cdots,G_{k+1}$ be the components of $G$ in order and let $g_i\ge0$ be the number of edges in $G_i$. Thus, $g_0=\mu$ and the other $g_i$ are some permutation of the ${\lambda_j}$ in $\lambda$. Let $d_0,\cdots,d_k$ be the missing edges in order. Thus $d_i$ is the edge between $G_i$ and $G_{i+1}$. For any $(G,e)\in \widetilde\cS_n(\mu;\lambda)$, let $\varphi(G,e)$ be given as follows.
\begin{enumerate}
\item $\varphi(G,e)=(G,e)$ if $e\notin G$.
\item $\varphi(G,e)=(G,d_0)$ if $e\in G_0\cup G_1$. Thus $\varphi^{-1}(G,d_0)$ has $\mu+g_1+1$ elements.
\item $\varphi(G,e)=(G,d_i)$ if $e\in G_{i+1}$ and $i\ge2$. So, $|\varphi^{-1}(G,d_i)|=g_{i+1}+1$ if $i\ge2$.
\end{enumerate}

Now we count the number of elements of $\widetilde\cS_n(\mu;\lambda)$ corresponding to each element of the right hand side of \eqref{eq: recursive count of S(mu,l)}.
\begin{enumerate}
	\item Each element of $\cS_n(\mu';\lambda')$ gives an element $(G,d_0)\in \widetilde\cS_n^p(\mu;\lambda)$ with $g_0=\mu,g_1=a$. This has $\mu'=\mu+a+1$ inverse image points in $\widetilde\cS_n(\mu;\lambda)$.
	\item Take $G''\in\cS_n(\mu;\lambda'')$ where $a=b$ and $c=2a+1$. There are $n_c''$ components of $G''$ of size $c$ (excluding $G_0''$). In each such component, the middle edge $p$ is removed giving two new components of $G=G''\backslash p$ both of size $a$. Then $(G,p)$ is the corresponding element of $\widetilde \cS_n^p(\mu;\lambda)$ with $a+1$ inverse image points in $\widetilde \cS_n^p(\mu;\lambda)$. This gives
	\[
	(a+1)n_c''=n_c''(c+1)X(a,a)
	\]
elements of $\widetilde \cS_n(\mu;\lambda)$ since $X(a,a)=\frac12$ and $c+1=2a+2$.
	\item Take $G''\in\cS_n(\mu;\lambda'')$ where $a\neq b$ and $c=a+b+1$. There are $n_c''$ components of $G''$ of size $c$ (excluding $G_0''$). In each such component there are two edges, say $p,q$, which can be deleted to produce an element of $\cS_n(\mu;\lambda)$. In $G''\backslash p$, the $L_c$ component becomes $L_a\coprod L_b$ and in $G''\backslash q$, the $L_c$ becomes $L_b\coprod L_a$. In the first case the inverse image in $\widetilde\cS_n(\mu;\lambda)$ has $b+1$ elements, in the second case it has $a+1$ elements for a total of
	\[
		n_c''(a+1+b+1)=n_c''(c+1)X(a,b)
	\]
	inverse image points since $X(a,b)=1$.
\end{enumerate}
Adding these up gives the lemma.
\end{proof}

\end{document}
